% Original Prism research-article skeleton for Nucleic Acids Research.
% Uses the unmodified OUP 1.5 class and its Modern Large page design.
% This is a writing starter, not an official NAR sample manuscript.
\documentclass[unnumsec,webpdf,modern,large]{oup-authoring-template}
\usepackage{booktabs,tabularx,siunitx}
\newcommand{\societylogo}{}
\sisetup{uncertainty-mode=separate}

\begin{document}
\journaltitle{Nucleic Acids Research}
\DOI{}
\copyrightyear{}
\pubyear{}
\access{Draft manuscript}
\appnotes{Research Article}
\firstpage{1}

\title[Short manuscript title]{Manuscript title}
\author[1,$\ast$]{First Author}
\author[1,2]{Second Author}
\address[1]{\orgdiv{Department}, \orgname{First institution},
  \orgaddress{\state{City}, \postcode{Postal code}, \country{Country}}}
\address[2]{\orgdiv{Department}, \orgname{Second institution},
  \orgaddress{\state{City}, \postcode{Postal code}, \country{Country}}}
\corresp[$\ast$]{Corresponding author.
  \href{mailto:first.author@example.org}{first.author@example.org}}

% One English paragraph, at most 200 words; no citations, figure references
% or URLs in a standard research-article abstract.
\abstract{Replace this paragraph with the biological question, approach,
verified findings and their significance. Use plain language appropriate to
the readership and define any necessary standard abbreviation. This
starter illustrates manuscript structure only. Its table values, equation
and graphic slots are synthetic and do not represent a study. Complete
the abstract and replace every instructional passage before submission.}
\maketitle
% Avoid production-only volume/issue placeholders and the draft title banner.
\thispagestyle{plain}
\pagestyle{plain}

\section{Introduction}\label{sec:introduction}
Introduce the biological question, the unresolved problem and the objective
of the study. Distinguish the contribution from relevant prior work.
The public citation here tests bibliography mechanics only; retain it only
if relevant to the manuscript~\cite{wilkinson2016fair}.

\section{Materials and Methods}\label{sec:methods}
\subsection{Study design and resources}
Describe the actual study design, source data, selection rules and
reagents or biological resources where applicable. Give identifiers and
versions that permit the work to be reproduced. Keep essential methods
in the main text; use the supplement for additional detail.

\subsection{Analysis and statistical methods}
Define the analysis, comparison conditions, independent repetitions,
sample sizes, uncertainty measures and statistical procedures actually
used. Record software versions and access to new code. Explain applicable
ethical approvals and data-access restrictions accurately.

Equation~\eqref{eq:residual} demonstrates mathematical cross-referencing.
Replace or remove this purely illustrative expression.
\begin{equation}
 E=\frac{1}{n}\sum_{i=1}^{n}(y_i-\widehat y_i)^2.
 \label{eq:residual}
\end{equation}

\section{Results}\label{sec:results}
Present verified observations in the order needed to answer the study
question. Figure~\ref{fig:study} and Table~\ref{tab:study} are synthetic
layout examples, not evidence.

\begin{figure}[t]
\centering
\fbox{\parbox[c][33mm][c]{0.87\linewidth}{\centering
Original study graphic\\[5pt]Replace this synthetic slot}}
\caption{Synthetic graphic slot. Describe the actual panels, units,
sample sizes and uncertainty when replacing it.}
\label{fig:study}
\par\smallskip\small\raggedright
\textbf{Alt text:} Replace with a concise description of the actual graphic.
\end{figure}

\begin{table}[t]
\centering\small
\caption{Synthetic numerical entries for typesetting only.}
\label{tab:study}
\begin{tabularx}{\linewidth}{@{}>{\raggedright\arraybackslash}X
 S[table-format=2.0] S[table-format=1.3(3)]@{}}
\toprule
Setting & {Count} & {Estimate $\pm$ SD}\\
\midrule
Example A & 24 & 0.320 +- 0.015\\
Example B & 36 & 0.415 +- 0.021\\
\bottomrule
\end{tabularx}
\par\smallskip\footnotesize\raggedright
All values are invented. Define the actual units and uncertainty measure
when replacing these entries.
\end{table}

\section{Discussion}\label{sec:discussion}
Interpret the findings in relation to the question, prior work and
biological relevance. Describe limitations and the evidence supporting
each conclusion; separate speculation from established results.
Refer to Supplementary Figure~S1, Supplementary Table~S1 and
Supplementary Equation~(S1) in the separate supplementary document
only where they support the final manuscript.

\section*{Data Availability}
Replace with the actual access routes, stable identifiers and versions
for all underlying data and code, together with any justified restrictions.
For newly reported software, provide the permanent archived code DOI
when required at revision or acceptance.

\section*{Supplementary Data}
Describe the supplied supplementary files and replace this instruction
with the journal's required availability statement when applicable.

\section*{Author Contributions}
Replace with the contributions agreed by every author, using appropriate
CRediT roles where applicable.

\section*{Acknowledgements}
Replace with confirmed acknowledgements, or remove if not applicable.

\section*{Funding}
Replace with confirmed funder names, grant identifiers and the source of
open-access funding where applicable. Do not infer a funding statement.

\section*{Conflict of Interest}
Replace with the authors' confirmed competing-interest disclosure.

\bibliographystyle{prism-nar}
\bibliography{references}
\end{document}
